For a positive integer $N$, an integer $\alpha$ is said to be an $\bemod$-th residue modulo $N$ if and only if there exists some integer $\beta$ such that $\alpha \equiv \beta^\bemod \pmod{N}$.
Let $\Z^\ast_N$ denote the set of integers that are relatively prime to $N$.
Let $\resmod{\bemod}{N}$ denote the set of $\bemod$-th residues modulo $N$ in $\Z^\ast_N$, and let $\nonresmod{\bemod}{N} = \Z^\ast_N \setminus \resmod{\bemod}{N}$.

Let $\GrpGen(\secparam, \bemod) \rightarrow (p, q, g)$ be a $\ppt$ algorithm such that (1) $p$ and $q$ are prime, (2) $\bemod$ divides $\eulertot(N)$, (3) $\bemod$ does not divide $\eulertot(N) / \bemod$, and (4) $g \in \nonresmod{\bemod}{N}$, where $N = pq$ and $\eulertot(N) = (p-1)(q-1)$ denotes Euler's totient function.
We assume there exists a bound $\bemodbound(\secpar) = 2^{\bigO{\secpar}}$ such that $p$ and $q$ output by $\GrpGen(\secparam, \bemod)$ satisfy $p < q < \bemodbound(\secpar)$.
The prime residuosity assumption informally says that for any prime $\bemod$, it is computationally hard to distinguish between the uniform distribution over $\resmod{\bemod}{N}$ and uniform distribution over $\Z^\ast_N$, where $N = pq$ is computed as the output of $\GrpGen$.

\begin{definition}[Prime Residuosity Assumption]\label{def:prime-residuosity}
	For all non-uniform polynomial time adversaries $\adv$, there exists a negligible function $\negl[]$ such that for all primes $\bemod$ and $\secpar \in \N$, we have
	\begin{equation*}
		\begin{gathered}
			\abs{\prob{\adv(\secparam, N, g, \alpha) = 1} - \prob{\adv(\secparam, N, g, \beta) = 1}} \\
			\le \negl,
		\end{gathered}
	\end{equation*}
	where $(p, g, q) \gets \GrpGen(\secparam, \bemod)$, $N = pq$, $\alpha \gets \Z^\ast_N$, and $\beta \gets \resmod{\bemod}{N}$.
\end{definition}

We also recall some useful facts about the structure of the group generated by $\GrpGen$.

\begin{lemma}[{\cite[Theorem 2.22]{THESIS:Ben87}}]\label{lemma:benaloh:unique-decomp}
	For all $\secpar \in \N$, all primes $\bemod \in \N$, every $(p, q, g)$ in the support of $\GrpGen(\secparam, \bemod)$ and all $\alpha \in \Z^\ast_N$, where $N = pq$, there exists a unique $x \in \zrange{0}{\bemod}$ and unique $\gamma \in \resmod{\bemod}{N}$ such that $\alpha \equiv g^x \gamma \bmod{N}$.
\end{lemma}

\begin{lemma}[{\cite[Theorem 2.24]{THESIS:Ben87}}]\label{lemma:benaloh:efficient-dlog}
	There exists an algorithm $\dlog$ such that for all $\secpar \in \N$, all polynomials $\poly[\cdot]$, all primes $\bemod \in \N$, every $(p, q, g)$ in the support of $\GrpGen(\secparam, \bemod)$ and all $\alpha \in \Z^\ast_N$, where $\bemod \le \poly$ and $N = pq$, $\dlog(p, q, g, \alpha)$ computes $x \in \zrange{0}{\bemod}$ in polynomial time in $\secpar$ such that $\alpha \equiv g^x \gamma \bmod{N}$ for some $\gamma \in \resmod{\bemod}{N}$.
\end{lemma}

